\chapter{2011:Dan Shechtman}

\label{chap:chemistry2011}

The Nobel Prize in Chemistry for the year 2011 was awarded to Dan Shechtman.

\textbf{Motivation:}

for the discovery of quasicrystals

\section{Dan Shechtman}

\subsection*{Early Life and Education}

Dan Shechtman was born in 1941 in Tel Aviv, Israel.

\subsection*{Career and Major Research}

Shechtman earned bachelor's, master's, and doctoral degrees in materials engineering from the Technion--Israel Institute of Technology in Haifa. In 1982, while on sabbatical at the National Bureau of Standards in the United States, he observed an electron diffraction pattern with fivefold symmetry from a rapidly cooled aluminum--manganese alloy, a symmetry forbidden for ordinary crystals.

\subsection*{Nobel Prize-winning Work}

The work for which Dan Shechtman was awarded the Nobel Prize in Chemistry in 2011 was described by the Nobel Committee as: "for the discovery of quasicrystals".

Shechtman's diffraction patterns showed a solid with long-range order but no translational periodicity, a phase that could not be described as a conventional crystal. The new material was named a quasicrystal, and his findings were published with colleagues in 1984.

\subsection*{Significance and Impact}

The discovery overturned the long-standing belief that all crystals must have repeating translational symmetry. Quasicrystals possess order, yet their atoms are arranged with rotational symmetries such as fivefold that crystals were thought to forbid.

\subsection*{Later Career and Honors}

Shechtman was professor of materials science at the Technion, where he is now professor emeritus, and he has held positions at research institutions in the United States and Europe.

\subsection*{Legacy}

Since his discovery, hundreds of quasicrystalline materials have been identified, and quasicrystals are now recognized as a distinct class of solid matter with applications in coatings and alloys.

\subsection*{Selected Publications}

"Metallic Phase with Long-Range Orientational Order and No Translational Symmetry" (Physical Review Letters, 1984)

\clearpage

\section*{Chapter Summary}

The 2011 prize recognized Dan Shechtman's discovery of quasicrystals, materials with ordered but non-periodic atomic arrangements. The finding changed the definition of a crystal and opened a new branch of materials science.
